% Source: round3/arithmetic/sharp_character_span.tex
% Reviewed branch fragment; requires theorem/proposition/proof environments.
\section{The exact energy floor of the full quadratic-character template}
The earlier constructions used one particular arrangement of a quadratic
character, constants, and the all-one polynomial. The following result
shows that their minimum energies cannot be improved by taking arbitrary
integer linear combinations of those same three ingredients, even with
independent shifts in the four channels. This is a restriction on a
construction template, not on unrestricted Hadamard existence.

Let $m>1$ be odd, $R_m=\Z[z]/(z^m-1)$, $z^*=z^{-1}$ and
$J=1+z+\cdots+z^{m-1}$. Suppose $L$ has coefficient zero at the
identity, all other coefficients in $\{\pm1\}$, and
\[
 LL^*=m\cdot1-J,\qquad L^*=\eta L,\qquad \eta\in\{+1,-1\}.
\]
For prime $m=p$, the quadratic character polynomial has these
properties with $\eta=(-1)^{(p-1)/2}$. Its coefficients take both signs:
augmentation of the norm identity forces $L(1)=0$.
Consider the four-channel family
\begin{equation}\label{s03:eq:fullspan}
 A_i=z^{r_i}(u_iL+v_i\cdot1+w_iJ),\qquad
 u_i,v_i,w_i,r_i\in\Z.
\end{equation}
Require every coefficient of every $A_i$ to be odd and
$\sum_iA_iA_i^*=4t\cdot1$.

\begin{theorem}[Sharp floor in the full character span]\label{s03:thm:spanfloor}
If $m\equiv3\pmod4$ and $L^*=-L$, the least possible energy in
\eqref{s03:eq:fullspan} is
\[
 T_-(m)=\frac{m^2+6m+21}{16}.
\]
If $m\equiv1\pmod4$, $m\ge5$, and $L^*=L$, the least possible energy is
\[
 T_+(m)=\frac{m^2+2m+5}{8}.
\]
Both minima are attained by explicit odd-coefficient choices. The
integers $T_\pm(m)$ are odd. These are global minima within the whole
family, not only within the earlier special constructions.
\end{theorem}
\begin{proof}
Shifts do not change the individual Hermitian norms and preserve
coefficient parity, so omit them. Since the nonidentity coefficients
of $L$ include both signs, coefficientwise oddness is equivalent to
\begin{equation}\label{s03:eq:spanparity}
 u_i\equiv v_i\equiv1-w_i\pmod2.
\end{equation}
Put $S=\sum_i u_i^2$, $V=\sum_i v_i^2$, and $W=\sum_i w_i^2$.
Since $LJ=0$, the norm sum is
\[
 (mS+V)\cdot1+
 \left(-S+2\sum_iw_iv_i+mW\right)J
 +(1+\eta)\left(\sum_i u_iv_i\right)L.
\]
The three polynomials $1,J,L$ are linearly independent over $\mathbb Q$.
Consequently scalarity requires
\begin{equation}\label{s03:eq:spanconditions}
 2w\mathbin{\cdot}v=S-mW,\qquad 4t=mS+V,
\end{equation}
and in the symmetric case also $u\mathbin{\cdot}v=0$.
The energy $t$ is odd, since $4m$ odd coefficients have squared sum
$4$ modulo eight. Also $W>0$: otherwise scalarity forces $S=0$,
contradicting~\eqref{s03:eq:spanparity}. Cauchy--Schwarz gives the useful
universal inequality
\begin{equation}\label{s03:eq:spanCS}
 t\ge\frac{mS+(S-mW)^2/(4W)}4
   =\frac{(S+mW)^2}{16W}.
\end{equation}

In the skew case, first suppose $W\ge2$. Then $t\ge m^2/8$, while
\[
 \frac{m^2}{8}-(T_-(m)-2)
 =\frac{(m-3)^2+2}{16}>0.
\]
Since $t$ and $T_-(m)$ are odd integers, this implies $t\ge T_-(m)$.
If $W=1$, exactly one $w_i$ is $+1$ or $-1$ and the other three
vanish. Thus three $u_i$ and three corresponding $v_i$ are odd,
giving $S\ge3$ and
\[
 V\ge\frac{(S-m)^2}{4}+3.
\]
It follows that
\[
 t\ge\frac{mS+(S-m)^2/4+3}{4}
 =\frac{(m+S)^2}{16}+\frac34
 \ge T_-(m).
\]
For attainment write $m=4k+3$, and take
\[
 (A_1,A_2,A_3,A_4)=(L+1,L+1,L+1,J-2k).
\]
The first three norm sums contribute $(3m+3)1-3J$, while the last
contributes $4k^2\cdot1+3J$, giving energy $T_-(m)=k^2+3k+3$.

In the symmetric case, $u\cdot v=0$ and
\eqref{s03:eq:spanparity} imply that the number of odd $w_i$ is even.
Hence $W$ is a positive even integer. If $W\ge4$,
\eqref{s03:eq:spanCS} gives $t\ge m^2/4$. Since
\[
 \frac{m^2}{4}-(T_+(m)-2)
 =\frac{(m-1)^2+10}{8}>0,
\]
odd integrality again gives $t\ge T_+(m)$. If $W=2$, exactly two
$w_i$ are $\pm1$; the other two vanish. Their corresponding $u_i,v_i$
are odd, so $S\ge2$. Applying Cauchy--Schwarz only on the two
nonzero-$w_i$ coordinates yields
\[
 V\ge\frac{(S-2m)^2}{8}+2.
\]
Thus
\[
 t\ge\frac{mS+(S-2m)^2/8+2}{4}
 =\frac{(S+2m)^2}{32}+\frac12\ge T_+(m).
\]
For attainment, write $m=4k+1$ and set
\[
 (A_1,A_2,A_3,A_4)=(L+1,L-1,J-2k,J-2k).
\]
The two $L$ norms have cancelling $L$ terms; their $J$ coefficient
is $-2$, while the last two norms have $J$ coefficient $+2$.
The scalar energy is $2k^2+2k+1=T_+(m)$.
\end{proof}

\begin{corollary}[Only three prime lengths can reach the binary boundary]
For the quadratic character polynomial at an odd prime $p$, the full
family~\eqref{s03:eq:fullspan} reaches energy $t=p$ if and only if
$p\in\{3,5,7\}$. In particular no prime $p\ge11$ can be settled
by rescaling, independently shifting, or taking arbitrary integer
linear combinations of a quadratic character, its zero-position
monomial, and $J$ in the four channels.
\end{corollary}
\begin{proof}
For skew lengths,
$T_-(p)-p=(p-3)(p-7)/16$; for symmetric lengths,
$T_+(p)-p=(p-1)(p-5)/8$. The displayed attaining constructions
reach the boundary at $3,5,7$ and have larger minimum energy at all
other odd primes. At the boundary every coefficient is a sign.
\end{proof}

This also covers affine relabelings of the quadratic character:
$\chi(aj+b)$ is, up to sign, a shift of $\chi(j)$, with its unique
zero shifted at the same time. It does not cover arbitrary changes
at additional positions, nonquadratic characters, or general sequence
quadruples. In particular it is not an obstruction to Hadamard matrices
of the corresponding orders.

\subsection{Fixing only two character channels already excludes infinitely many primes}
Let $p\equiv3\pmod4$, and let two channels be the binary sequences
$z^r(L+\varepsilon)$ and $z^s(L+\delta)$, where
$\varepsilon,\delta\in\{\pm1\}$. Their row sums have squares one.
If two otherwise unrestricted binary channels complete these to a
four-sequence complementary set, augmentation forces their row sums
$a,b$ to satisfy
\[
 a^2+b^2=4p-2.
\]
Both are odd; equivalently
$((a+b)/2)^2+((a-b)/2)^2=2p-1$.
For $p\equiv11\pmod{36}$, the integer $2p-1$ is three modulo nine,
so its $3$-adic valuation is exactly one. A sum of two squares cannot
have that valuation: divisibility by three forces both summands to
be divisible by three, making the valuation at least two. Thus no
such completion exists at these primes, regardless of how the other
two channels are chosen. Dirichlet's theorem supplies infinitely many
primes in the class $11\pmod{36}$; an English translation of the
original memoir is available at \url{https://arxiv.org/abs/0808.1408}.

Fixing three binary character channels imposes the stronger necessary
condition that $4p-3$ be a square. With only one fixed channel, the
remaining augmentation equation is a sum of three odd squares equal
to $4p-1\equiv3\pmod8$ and has no obstruction from the three-square
theorem. This does not assert that the full sequence constraints are
then solvable.

\paragraph{Scope and reproduction.}
The exact checker constructs the attaining matrices for primes through
$101$, checks their periodic autocorrelations, samples scalar templates
with arbitrary integer coefficients, and verifies the two-channel
congruence obstruction at small primes. The infinite assertions rest
on the proofs, not on these tests. These restrictions identify the
precise need to leave the quadratic-character affine span; they do not
improve the main Hadamard exponent bound or prove new all-order existence.

